\section{Cross-validation for hyperparameter selection}
\label{sec:cross_validation}

To obtain a robust surrogate model while avoiding overfitting to a single train--validation split, cross-validation (CV) was used for hyperparameter selection in both stages of the learning framework. Since the dataset is highly imbalanced with respect to the observed gap-index category, the CV procedure was designed separately for the gap-index classification task and the gap-edge regression task. The classification model was tuned using stratified CV, whereas the regression models were tuned within each gap-index branch.

\subsection{Cross-validation for gap-index classification}

For the first-stage gap-index classifier, a stratified five-fold CV was performed on a randomly selected $20\%$ subset of the full dataset. The sampling procedure preserved the class proportions of the original dataset, and each fold was constructed such that the four classes, namely no gap, gap index 2, gap index 3, and other gap indices, remained represented in both the training and validation partitions. This is important because the gap-index distribution is strongly skewed toward gap index 3, while the no-gap and other-index classes contain far fewer samples.

The classifier was evaluated primarily using the macro-averaged F1 score,
\begin{equation}
    F_{1,\mathrm{macro}}
    =
    \frac{1}{C}
    \sum_{c=1}^{C}
    F_{1}^{(c)},
\end{equation}
where $C=4$ is the number of classes. Compared with the overall accuracy, the macro-F1 score gives equal weight to each class and is therefore more suitable for imbalanced multi-class classification. Balanced accuracy and per-class F1 scores were also recorded to monitor the minority classes.

Several configurations were examined by varying the focal-loss parameter, label smoothing, weighted sampling strength, dropout rate, learning rate, network width, and hidden dimension. Table~\ref{tab:classifier_cv} summarizes representative CV results. The best configuration was obtained using focal loss with $\gamma=1.0$ and a moderate weighted-sampling power of $0.25$, which achieved the highest mean macro-F1 score.

\begin{table}[htbp]
    \centering
    \caption{Five-fold cross-validation results for the gap-index classifier on a stratified $20\%$ subset of the full dataset.}
    \label{tab:classifier_cv}
    \begin{tabular}{lcccccc}
        \toprule
        Configuration & Macro-F1 & Balanced acc. & Accuracy & No gap F1 & Gap 2 F1 & Gap 3 F1 \\
        \midrule
        $\gamma=1.0$, sampler $0.25$ & 0.7681 & 0.8522 & 0.9599 & 0.6350 & 0.8022 & 0.9805 \\
        Label smoothing $0.05$       & 0.7675 & 0.8602 & 0.9568 & 0.6550 & 0.7891 & 0.9785 \\
        Hidden width $256$           & 0.7637 & 0.8604 & 0.9557 & 0.6445 & 0.7852 & 0.9779 \\
        Dropout $0.05$               & 0.7626 & 0.8627 & 0.9558 & 0.6444 & 0.7892 & 0.9779 \\
        Learning rate $2\times10^{-4}$ & 0.7595 & 0.8642 & 0.9567 & 0.6443 & 0.7950 & 0.9785 \\
        \bottomrule
    \end{tabular}
\end{table}

\subsection{Cross-validation for gap-edge regression}

After the gap-index classifier determines the gap category, the lower and upper gap-edge frequencies are predicted by branch-specific regression models. Therefore, the regression task was treated independently for each gap-index branch rather than as a single mixed regression problem. In this work, the gap-index-3 branch was selected for detailed hyperparameter tuning because it contains the largest number of valid samples and is the most statistically supported branch of the dataset.

For the gap-index-3 regressor, five-fold CV was performed using a random $20\%$ subset of all gap-index-3 samples. This corresponds to $319{,}333$ samples selected from $1{,}596{,}663$ available gap-index-3 geometries. The selected samples were split into five folds, with approximately $63{,}866$ samples used for validation in each fold. For each candidate configuration, the model was trained for five epochs on four folds and evaluated on the remaining fold. No sample was shared between the training and validation partitions within a fold.

The regression target is the pair of lower and upper complete band-gap edge frequencies,
\begin{equation}
    \mathbf{y}
    =
    \left[
    f_{\mathrm{low}},
    f_{\mathrm{up}}
    \right],
\end{equation}
and the primary selection metric was the mean absolute error averaged over the two band-edge frequencies,
\begin{equation}
    \mathrm{MAE}_{\mathrm{edge}}
    =
    \frac{1}{2}
    \left(
    \mathrm{MAE}_{\mathrm{low}}
    +
    \mathrm{MAE}_{\mathrm{up}}
    \right).
\end{equation}
The gap-width error was also computed as a derived metric, where the predicted width is given by $f_{\mathrm{up}}-f_{\mathrm{low}}$.

The candidate configurations covered changes in learning rate, dropout rate, hidden dimension, network depth, and the smooth-$L_1$ loss parameter. The CV results are summarized in Table~\ref{tab:gap3_regressor_cv}. The best configuration used a learning rate of $5\times10^{-4}$ with the baseline architecture, achieving an average band-edge MAE of $0.1649~\mathrm{kHz}$ across the five folds.

\begin{table}[htbp]
    \centering
    \caption{Five-fold cross-validation results for the gap-index-3 lower/upper band-edge regressor using $20\%$ of all gap-index-3 samples. The selection metric is the mean MAE of the lower and upper gap-edge frequencies.}
    \label{tab:gap3_regressor_cv}
    \begin{tabular}{lccccc}
        \toprule
        Configuration & Edge MAE (kHz) & Lower MAE (kHz) & Upper MAE (kHz) & Width MAE (kHz) & $R^2_{\mathrm{up}}$ \\
        \midrule
        Learning rate $5\times10^{-4}$ & 0.1649 & 0.0781 & 0.2516 & 0.2605 & 0.9975 \\
        Learning rate $4\times10^{-4}$ & 0.1764 & 0.0817 & 0.2711 & 0.2809 & 0.9971 \\
        Dropout $0.05$                 & 0.1857 & 0.0830 & 0.2885 & 0.2927 & 0.9967 \\
        Smooth-$L_1$ $\beta=0.5$       & 0.1899 & 0.0845 & 0.2953 & 0.3017 & 0.9965 \\
        Six message-passing layers     & 0.1981 & 0.0885 & 0.3077 & 0.3173 & 0.9963 \\
        Hidden width $256$             & 0.1987 & 0.1011 & 0.2962 & 0.3110 & 0.9965 \\
        Baseline                       & 0.1989 & 0.0913 & 0.3064 & 0.3186 & 0.9963 \\
        Ten message-passing layers     & 0.1994 & 0.0884 & 0.3105 & 0.3178 & 0.9963 \\
        Dropout $0.15$                 & 0.2063 & 0.0950 & 0.3175 & 0.3260 & 0.9961 \\
        Hidden width $192$             & 0.2113 & 0.0963 & 0.3262 & 0.3353 & 0.9958 \\
        Dropout $0.20$                 & 0.2169 & 0.1009 & 0.3329 & 0.3402 & 0.9958 \\
        Learning rate $2\times10^{-4}$ & 0.2226 & 0.0974 & 0.3477 & 0.3566 & 0.9952 \\
        \bottomrule
    \end{tabular}
\end{table}

Based on these results, the selected hyperparameters for the gap-index-3 regressor were hidden dimension $224$, eight message-passing layers, dropout rate $0.10$, weight decay $10^{-5}$, smooth-$L_1$ loss with $\beta=1.0$, and learning rate $5\times10^{-4}$. The CV results show that the learning rate has a stronger effect on the regression accuracy than moderate changes in network width or depth. The same CV-based selection strategy can be applied to the gap-index-2 branch before final training if an independently optimized gap-index-2 expert is required.
